\section{Prime supply, collisions, and the arithmetic grid}
\label{app:arithmetic-details}

The accepting density has been defined on complete cell words before
actual primes are assigned. We first prove the collision estimate for
an arbitrary such weight, then show that every supported occupancy has
sufficient squarefree harmonic mass. The final comparison keeps the
physical optimizer fixed throughout the passage to complete prime cells.

\subsection{A fractional collision kernel for correlated label weights}
\label{sec:correlated-collision-kernel}
We prove a prime-cell majorant for arbitrary nonnegative label weights,
including weights coupling different slots. The triangularization is
related to the unweighted argument of \cite[Section~7.1]{Kouk}; we give
the finite history construction and the complete weighted argument here.

\subsubsection{Cells, slots, and the arithmetic moment}
Fix $k\ge1$, put $b=\log2$, and set
\[
 \eta_i=\frac{\log(k-i+2)}{k-i+1},\qquad \rho_i=e^{\eta_i},
 \qquad y_0=3\le y_1\le\cdots\le y_k.
\]
Use pairwise disjoint prime cells
$D_{i,t}=(\lambda_{i,t-1},\lambda_{i,t}]\cap\mathbb P$
for $1\le t\le v_i$, satisfying
\begin{align}
 &\lambda_{i,t-1}\ge3,\qquad
   \sum_{p\in D_{i,t}}p^{-1}\le\eta_i,\notag\\
 &\log\lambda_{i,t-1}\ge c_*\log y_{i-1}\rho_i^t,\qquad
 \left|\log\frac{\log y_i}{\log y_{i-1}}-\eta_iv_i\right|
       \le A_* .                                      \label{eq:cc-cell}
\end{align}
The constants $c_*>0,A_*<\infty$ are fixed. The elementary greedy
cells described below have $c_*^{-1},A_*\ll_k1$. Empty blocks with
$v_i=0$ are allowed; the last inequality still applies to them.

Fix integers $g_{i,t}\ge0$, and define
\[
 r_i=\sum_tg_{i,t},\quad R_i=\sum_{h\le i}r_h,\quad R_0=0,
 \quad G_g=\prod_{i,t}g_{i,t}!,\quad M_g=\prod_i\eta_i^{r_i}.
\]
The slots in block $i$ are $\mathcal B_i=(R_{i-1},R_i]\cap\mathbb Z$.
Assign its first $g_{i,1}$ slots to cell $1$, its next $g_{i,2}$
to cell $2$, and so on. Write $E_g(q)$ for the assigned cell of a
positive slot $q$, and set $E_g(0)=0$. No ordering of prime values
inside one cell is imposed.

A label is a map $Y$ assigning to a slot $q\in\mathcal B_i$ one
color in $C_q=\{0,i,i+1,\ldots,k\}$. Equivalently,
\[
 \mathcal P_{\mathbf r}=
 \{(Y_1,\ldots,Y_k):Y_j\subseteq\{1,\ldots,R_j\},\quad
                  Y_j\cap Y_h=\varnothing\ (j\ne h)\},
\]
where unused slots have color zero. Let
$F:\mathcal P_{\mathbf r}\to[0,\infty)$ be finite, independent
of all prime values, and invariant under every permutation of slots
within each assigned cell. No product structure is assumed for $F$.
It may depend on the entire occupancy vector $g$.

Let $\mathcal A(g)$ consist of squarefree tuples
$\mathbf a=(a_1,\ldots,a_k)$ having $g_{i,t}$ prime factors from
each $D_{i,t}$, with all prime factors coming from these cells.
Choose any named-slot representation $\mathbf p$ of such a tuple.
For labels $Y,Z$, put
\[
 Y\sim_{\mathbf p}Z
 \quad\Longleftrightarrow\quad
 \left|\sum_{q\in Y_j}\log p_q-\sum_{q\in Z_j}\log p_q\right|<b
                     \quad(1\le j\le k).
\]
For $1<P\le2$, define the arithmetic moment and its harmonic sum by
\begin{align}
 W_F^P(\mathbf a)
   &=\sum_{Y\in\mathcal P_{\mathbf r}}F(Y)
           \left(\sum_{Z:\,Y\sim_{\mathbf p}Z}F(Z)\right)^{P-1},
       \notag\\
 \mathcal W_F^P(g)
   &=\sum_{\mathbf a\in\mathcal A(g)}
                  \frac{W_F^P(\mathbf a)}{a_1\cdots a_k}.
                                                  \label{eq:cc-moment}
\end{align}
Within-cell invariance makes these definitions independent of the
chosen representation. Distinct labels correspond to distinct divisor
tuples because all the primes are distinct.

\begin{lemma}[Elementary cell supply]\label{lem:cc-cell-supply}
For every ordered cutoff vector above there are cells containing
exactly the primes $\max\{k,y_{i-1}\}<p\le y_i$ and satisfying
\eqref{eq:cc-cell} with constants depending only on $k$.
\end{lemma}
\begin{proof}
For one block put $u=\max\{k,y_{i-1}\}$ and $\eta=\eta_i$.
First greedily partition all primes above $u$ into successive maximal
finite cells whose reciprocal mass is at most $\eta$. This is
possible because every such prime satisfies $p^{-1}<\eta$ and the
prime reciprocal sum diverges. Let $\lambda_t$ be the last prime
in complete cell $t$, with $\lambda_0=u$, and let $q_{t+1}$ be the
first prime in the next cell. The mass $H_t$ of complete cell $t$
satisfies
\[
       \eta-q_{t+1}^{-1}<H_t\le\eta.
\]
All but at most $Q_k=\lceil2/\min_i\eta_i\rceil$ cells have
$q_{t+1}>Q_k$, and therefore $H_t\ge\eta/2$ for those cells.
Mertens's estimate, uniformly for $3\le u\le v$, gives
\[
 \sum_{u<p\le v}\frac1p
       =\log\frac{\log v}{\log u}+O(1).
\]
It follows that
$\log\lambda_t\ge c_k\log u\,e^{\eta t/2}$.
Consequently
\[
 \sum_{t\ge1}q_{t+1}^{-1}
 \le\sum_{t\ge1}\exp(-c_k\log3\,e^{\eta t/2})\ll_k1.
\]
Summing the deficits in $H_t$ now gives
$\sum_{s\le t}H_s=t\eta+O_k(1)$, whence
\[
           \log\lambda_t=\log u\,e^{t\eta+O_k(1)}.
\]
Truncate this infinite sequence at $y_i$. Retain each nonempty
cell intersecting $(u,y_i]$ and replace the endpoint of the last
retained cell by $y_i$. There is at most one incomplete mass,
whose additional deficit is at most $\eta$. Thus the same cumulative
mass argument gives $\log(\log y_i/\log u)=\eta v_i+O_k(1)$
when $y_i>u$, including the case of no retained primes.
If $y_i\le u$, use no cells and note that
$\log y_i/\log y_{i-1}$ is bounded above and below in terms of $k$.
Finally $\log u/\log y_{i-1}\asymp_k1$; the estimates at preceding
endpoints, including $t=1$, give \eqref{eq:cc-cell}.
\end{proof}

\subsubsection{All histories and their explicit partitions}
A history $h=(\sigma,\mathbf I)$ consists of a permutation
$\sigma=(\sigma_1,\ldots,\sigma_k)$ of $\{1,\ldots,k\}$ and
integers satisfying
\begin{equation}\label{eq:cc-history-support}
       0\le I_j\le R_{\sigma_j},\qquad
       \text{all positive }I_j\text{ are distinct}.
\end{equation}
Assign a positive pivot $I_j$ to its slot block and every zero pivot
to block $1$. Define
\[
 \mathcal J_i(h)=\{j:I_j\text{ is assigned to block }i\},\quad
 n_i(h)=|\mathcal J_i(h)|,\quad J_i(h)=\sum_{s\le i}n_s(h).
\]
Let $\mathfrak H_{\mathbf r}$ be the histories in
\eqref{eq:cc-history-support} with $J_i(h)\ge i$ for every $i$.
Necessarily $J_k(h)=k$. No monotonicity of $I_1,\ldots,I_k$ is
imposed. Zero pivots can precede positive pivots.

Put $A_j(h)=\{\sigma_j,\ldots,\sigma_k\}$. For a slot $q$,
partition its color set $C_q$ by the following explicit equivalence:
\begin{equation}\label{eq:cc-slot-partition}
 c\equiv_{h,q}c'\quad\Longleftrightarrow\quad
 1_{c\in A_j(h)}=1_{c'\in A_j(h)}
       \ \text{for every }j\text{ with }I_j<q.
\end{equation}
Color $0$ belongs to none of the $A_j(h)$. Write $\Pi_{h,q}$ for
the nonempty equivalence cells. Declare $Z$ compatible with $Y$,
written $Z\equiv_hY$, if their colors lie in the same cell of
$\Pi_{h,q}$ for every slot $q$. Equivalently,
\begin{equation}\label{eq:cc-union-compatibility}
 \bigcup_{s=j}^k\bigl(Z_{\sigma_s}\cap(I_j,R_k]\bigr)
 =\bigcup_{s=j}^k\bigl(Y_{\sigma_s}\cap(I_j,R_k]\bigr)
                         \quad(1\le j\le k).
\end{equation}
This equivalence follows simply by checking the membership of each
individual slot in the two unions. It includes repeated zero pivots,
empty blocks, and arbitrary ties among cell indices.

Define
\begin{align}
 \Phi_g(h)&=\prod_{i=1}^k
  \rho_i^{-(k-J_i(h))v_i-\sum_{j\in\mathcal J_i(h)}E_g(I_j)},
                                                     \label{eq:cc-phi}\\
 M_F(Y;h)&=\sum_{Z:\,Z\equiv_hY}F(Z),\qquad
 \Psi_{h,P}(F)=\sum_YF(Y)M_F(Y;h)^{P-1}.
                                                     \label{eq:cc-psi}
\end{align}
All these are explicit finite quantities. The enlarged history set
can include histories not generated by any pair; their nonnegative
contributions cause no problem in an upper bound.

\begin{theorem}[Correlated fractional collision majorant]
\label{thm:correlated-collision-kernel}
Under \eqref{eq:cc-cell}, uniformly for $1<P\le2$ and every weight
$F$ specified above,
\begin{equation}\label{eq:cc-kernel}
 \boxed{\displaystyle
 \mathcal W_F^P(g)
 \le C(k,c_*,A_*)\,
       \frac{\prod_i\eta_i^{r_i}}{\prod_{i,t}g_{i,t}!}
       \sum_{h\in\mathfrak H_{\mathbf r}}
                 \Phi_g(h)^{P-1}\Psi_{h,P}(F).}
\end{equation}
For the cells of Lemma~\ref{lem:cc-cell-supply} the constant is
$C_k$ only. It is independent of $P,\mathbf y,\mathbf r,g,F$.
\end{theorem}

\begin{proof}
Every arithmetic tuple in $\mathcal A(g)$ has exactly $G_g$ named
representations by distinct primes in the prescribed cells. Hence,
writing $d\mu(\mathbf p)=(\prod_qp_q)^{-1}$ for this finite
named-prime measure and $\delta=P-1\in(0,1]$, we have the exact
identity
\begin{equation}\label{eq:cc-factorial-identity}
 \mathcal W_F^P(g)=\frac1{G_g}
   \sum_Y F(Y)\sum_{\mathbf p\text{ distinct}}d\mu(\mathbf p)
       \left(\sum_ZF(Z)1_{Y\sim_{\mathbf p}Z}\right)^\delta.
\end{equation}
In particular the factorial is present before any prime is deleted.
The total mass of this measure is at most $M_g$ by
\eqref{eq:cc-cell}. For any nonnegative function $f$, finite-measure
H\"older (or concavity) gives
\begin{equation}\label{eq:cc-prime-interpolation}
        \sum d\mu\,f^\delta
        \le M_g^{1-\delta}\left(\sum d\mu\,f\right)^\delta.
\end{equation}
This includes $\delta=1$; if the measure has zero mass the desired
theorem is immediate.

We next prove a linear prime bound for a fixed pair $Y,Z$. Its
canonical history is defined without reference to any prime values.
Starting with all coordinate colors remaining, at step $j$ let
$U_j$ be the set of slots belonging to exactly one of
$Y_c\triangle Z_c$, as $c$ runs through the remaining coordinates.
If $U_j\ne\varnothing$, choose $I_j=\max U_j$, let $\sigma_j$
be the unique remaining coordinate containing that slot in its
symmetric difference, and remove $\sigma_j$. If
$U_j=\varnothing$, choose $I_j=0$, let $\sigma_j$ be the largest
remaining coordinate, and remove it. Continue for exactly $k$ steps.

At a positive step the pivot lies in
$Y_{\sigma_j}\triangle Z_{\sigma_j}$, so $I_j\le R_{\sigma_j}$.
It is absent from every later remaining difference, and therefore
cannot be selected again. A zero step need not be followed only by
zero steps: deleting its coordinate can create new exclusive slots.
Every removed coordinate $c\le i$ supplies a pivot in the first
$i$ blocks, where zeros are counted in block $1$. Thus $J_i\ge i$,
and the canonical history belongs to $\mathfrak H_{\mathbf r}$.

For disjoint-coordinate labels the exclusive set just used equals
the symmetric difference of the unions of the remaining coordinate
sets. Indeed a slot has at most one first-label color and at most
one second-label color: if both colors remain and are different it
occurs in two differences, and if exactly one remains it occurs in
one. Maximality of $I_j$, including $I_j=0$, therefore proves
\eqref{eq:cc-union-compatibility} for this canonical history.

Fix all nonpivot primes. Retain the collision row for each positive
pivot step and order these rows and their pivot columns by step.
The coefficient of a previously chosen pivot in any later row is
zero, and each diagonal coefficient is $1$ or $-1$. The resulting
square matrix is upper triangular, with all entries in
$\{0,1,-1\}$ and dimension at most $k$. Its inverse has row sums
bounded by $2^k$: this follows either by back substitution or the
finite expansion of the inverse of a unit triangular matrix.
Discarding the zero-step rows only enlarges the feasible set.
It follows that, for the fixed nonpivot primes, all positive pivot
logarithms lie in a product of intervals of length at most
$B_k=2^{k+1}b$. The interval centers depend only on the nonpivot
primes and the labels. This is containment in a rectangle, not an
independence assertion about collision constraints.

Chebyshev's prime-counting upper bound implies, for a fixed positive
$B_k$, that
\begin{align}
 \sum_{p\in D_{i,t}\cap[X,e^{B_k}X]}\frac1p
 &\le\frac{C_k}{\log\max\{3,\lambda_{i,t-1},X\}}
 \le\frac{C_kc_*^{-1}\rho_i^{-t}}{\log y_{i-1}}.
                                                     \label{eq:cc-pivot-mass}
\end{align}
For the first inequality, if the intersection is nonempty take
$L=\max\{\lambda_{i,t-1},X\}\ge3$ and contain it in
$[L,e^{B_k}L]$; then bound its number of primes by
$\pi(e^{B_k}L)$ and each reciprocal by $L^{-1}$.
After integrating the rectangle, discard all distinctness conditions
and use the mass bound $\eta_i$ for every nonpivot slot. At most
$k$ factors $\eta_i^{-1}$ are lost, a constant depending only on
$k$. Thus, for the canonical history $h(Y,Z)$,
\begin{equation}\label{eq:cc-physical-kernel}
 \sum_{\mathbf p\text{ distinct}}d\mu(\mathbf p)
               1_{Y\sim_{\mathbf p}Z}
 \le C(k,c_*)M_g
     \prod_i(\log y_{i-1})^{-n_i(h)}
     \rho_i^{-\sum_{j\in\mathcal J_i(h)}E_g(I_j)}.
\end{equation}
Here a zero pivot deletes no prime. We inserted a factor
$(\log y_0)^{-1}$ for each such pivot; their number is at most
$k$, so this only changes the displayed constant.

The physical cutoff factors equal the factors in
\eqref{eq:cc-phi} up to the stated constant. More explicitly,
summation by parts gives
\begin{align*}
 \sum_i n_i(h)\log\log y_{i-1}
 &=k\log\log3+
   \sum_{i<k}(k-J_i(h))\log\frac{\log y_i}{\log y_{i-1}}\\
 &=\sum_i(k-J_i(h))\eta_iv_i+O(k+k^2A_*).
\end{align*}
Therefore the left side of \eqref{eq:cc-physical-kernel} is at
most $C(k,c_*,A_*)M_g\Phi_g(h(Y,Z))$.

For fixed $Y$, partition the competitors $Z$ by their canonical
histories, and enlarge each generating class to all $Z\equiv_hY$.
Weights are nonnegative, so
\[
 \sum_ZF(Z)\sum_{\mathbf p\text{ distinct}}d\mu(\mathbf p)
                         1_{Y\sim_{\mathbf p}Z}
 \le C(k,c_*,A_*)M_g
            \sum_{h\in\mathfrak H_{\mathbf r}}\Phi_g(h)M_F(Y;h).
\]
Apply \eqref{eq:cc-prime-interpolation} in
\eqref{eq:cc-factorial-identity}, and use
$(\sum_h a_h)^\delta\le\sum_h a_h^\delta$ for $a_h\ge0$.
The powers of $M_g$ combine exactly to $M_g$.
Taking the preceding constant at least one, its $\delta$ power is
bounded by the same constant uniformly for $0<\delta\le1$.
Finally sum against $F(Y)$ to obtain \eqref{eq:cc-kernel}.
No factor $g_{i,t}$ or falling factorial is introduced when pivot
slots are removed: their choices already occur in the history sum,
and $G_g$ is the exact original multiplicity.
\end{proof}

\subsubsection{The correlated partition term}
Let $\mathcal Q_h$ be the partition of the whole label space into
the equivalence classes of $\equiv_h$. This is the Cartesian
partition induced by the slot partitions \eqref{eq:cc-slot-partition},
but the weights of its classes need not factor. For every
nonnegative label weight, the exact identity is
\begin{equation}\label{eq:cc-global-power-sum}
 \Psi_{h,P}(F)
       =\sum_{Q\in\mathcal Q_h}\left(\sum_{Y\in Q}F(Y)\right)^P.
\end{equation}
Indeed $M_F(Y;h)$ is constant on $Q$, with value
$\sum_{Z\in Q}F(Z)$, so the anchors in $Q$ supply the remaining
first power. If $F(Y)=\prod_qw_{q,Y(q)}$, and only in that
additional case for this factorization, the right side becomes
\[
        \prod_q\sum_{B\in\Pi_{h,q}}
                          \left(\sum_{c\in B}w_{q,c}\right)^P.
\]
For a Gibbs weight depending on aggregate counts from several slots,
\eqref{eq:cc-global-power-sum} must be retained or estimated with
its correlations intact. Substituting its one-slot marginals in the
last product is not justified.

The theorem is pointwise in $g$. Thus one may multiply both sides
by any nonnegative occupancy test $T(g)$ and sum over $g$, allowing
$F=F_g$ separately at every occupancy. This preserves exactly the
nominal harmonic factorial mass $M_g/G_g$ and the same constant.
It does not assert a lower bound for the actual prime mass, nor
identify the occupancy distribution conditioned on a collision.
In particular, a bound for the weighted history and occupancy sum
on the right of \eqref{eq:cc-kernel}, with all required polynomial
factors, is still needed to deduce a sharp lower bound for $K_k$.

\subsection{Squarefree mass at every supported occupancy}
\label{arithdetail:supply}
For a nonempty finite prime cell $D$, put
\[
 H_D=\sum_{p\in D}p^{-1},\quad p_D=\min D,\quad
 e_g(D)=\sum_{\substack{p_1<\cdots<p_g\\p_j\in D}}(p_1\cdots p_g)^{-1},
 \qquad e_0(D)=1.
\]
Drawing $g$ independent primes with probabilities $1/(pH_D)$ gives
the exact identity
\[
 g!e_g(D)=H_D^g\mathbb P\{\text{all draws distinct}\}.
\]
Conditional on $s$ distinct previous draws, their total probability
mass is at most $s/(H_Dp_D)$. If $(g-1)/(H_Dp_D)\leq1/2$, then
\[
 \mathbb P\{\text{all distinct}\}
 \geq\prod_{s=0}^{g-1}\left(1-\frac{s}{H_Dp_D}\right)
 \geq\exp\left(-\frac{g(g-1)}{H_Dp_D}\right).
\]
For a complete greedy cell with target mass $\eta$,
$0\leq\eta-H_D<1/p_{\rm next}\leq1/p_D$.
If $p_D\geq2/\eta$ and $(g-1)/(H_Dp_D)\leq1/2$, then
$H_D\geq\eta/2$ and
\begin{equation}
 1\geq\frac{g!e_g(D)}{\eta^g}
 \geq\exp\left(-\frac{2g^2}{\eta p_D}\right).
 \label{arithdetail:single-supply}
\end{equation}
This includes $g=0$.

Use only complete retained cells $D_{i,j}$, omitting the final
incomplete cell, and set $p_{i,j}=\min D_{i,j}$,
$\eta_*=\min_i\eta_i$, and
\[
 W_g=\prod_{i,j}\frac{\eta_i^{g_{i,j}}}{g_{i,j}!},\qquad
 S_g^{\rm all}=\prod_{i,j}e_{g_{i,j}}(D_{i,j}).
\]
The second expression is exactly the squarefree harmonic tuple mass,
since the cells are disjoint. Provided the single-cell capacity
conditions above hold, multiplication of
\eqref{arithdetail:single-supply} gives
\begin{equation}
 1\geq S_g^{\rm all}/W_g
 \geq\exp\left(-\frac2{\eta_*}\sum_{i,j}\frac{g_{i,j}^2}{p_{i,j}}\right).
 \label{arithdetail:product-supply}
\end{equation}

We verify these conditions uniformly on the accepting support.
Let $L_i=\eta_iv_i$, suppose the total count obeys $r_i/L_i\leq D_0$,
and each of its $h_i$ fine colors satisfies at cell boundaries
\[
 \left|N_{i,c}(t)-\frac{t}{L_i}q_{i,c}\right|
 \leq H_0+\theta\min(t,L_i-t).
\]
An empty block has count zero. Summing over colors gives, for the
uncolored cumulative count $G_i(j)=\sum_{s\leq j}g_{i,s}$,
\[
 \left|G_i(j)-\frac j{v_i}r_i\right|
 \leq h_iH_0+h_i\theta\eta_i\min(j,v_i-j).
\]
Using $g_{i,j}\leq G_i(j)$ and
$g_{i,j}\leq r_i-G_i(j-1)$ separately, and putting
$A_0=(k+1)H_0$, $A_1=(\max_i\eta_i)(D_0+(k+1)\theta)$, proves
\begin{equation}
 g_{i,j}\leq A_0+A_1\min(j,v_i-j+1).
 \label{arithdetail:two-sided-tube}
\end{equation}
Only the existence of one accepted labeling is needed. The uncolored
occupancy is the same for every labeling, and recoloring changes none
of it. The named REG bounds in the construction imply the displayed
linear tubes since $t^\zeta\leq1+t$ for $0<\zeta<1$.

Choose all retained primes greater than a fixed $Q\geq2/\eta_*$.
Each complete cell has mass at least $\eta_i/2$. The previous $j-1$
retained cells and the reciprocal-integer bound imply
\[
 p_{i,j}\geq\frac Q2e^{(j-1)\eta_i/2}.
\]
Let $q=e^{-\eta_*/2}$ and define
\begin{align*}
 B_1&=\frac{A_0}{1-q}+\frac{A_1}{(1-q)^2},\\
 B_2&=\frac{A_0^2}{1-q}
   +\frac{2A_0A_1}{(1-q)^2}
   +\frac{A_1^2(1+q)}{(1-q)^3}.
\end{align*}
These are the sums of $(A_0+A_1j)q^{j-1}$ and its squared-polynomial
counterpart. Thus
\[
 \frac{g_{i,j}}{H_{i,j}p_{i,j}}\leq\frac{4B_1}{\eta_*Q},
 \qquad
 \sum_{i,j}\frac{g_{i,j}^2}{p_{i,j}}\leq\frac{2kB_2}{Q}.
\]
Taking $Q\geq\max(3,2/\eta_*,8B_1/\eta_*)$ makes every capacity
condition valid, and yields pointwise
\begin{equation}
 1\geq S_g^{\rm all}/W_g
 \geq\exp(-4kB_2/(\eta_*Q))>0.
 \label{arithdetail:uniform-supply}
\end{equation}
This is a multiplicative bound at every supported occupancy; no
absolute probability of repeated primes is subtracted from a rare mass.

The same support has a fixed size bound. Write $P_{i,j}=\max D_{i,j}$.
The $v_i-j$ later cells and the uniform Mertens difference estimate give
\[
 \log P_{i,j}\leq e^{C_{\rm M}}\log y_i\,
                         e^{-\eta_i(v_i-j)/2}.
\]
The right-hand tube in \eqref{arithdetail:two-sided-tube} therefore implies
\[
 \log a_i\leq\sum_jg_{i,j}\log P_{i,j}
 \leq e^{C_{\rm M}}B_1\log y_i.
\]
Choose the allowed size exponent at least
$A_{\rm tube}=\max(1,e^{C_{\rm M}}B_1)$.
Then the restricted actual mass equals $S_g^{\rm all}$ on every
supported occupancy, proving \eqref{lower:actual-supply}. The support
constants are fixed before this roughness and grid choice.

\subsection{Cancellation and one arithmetic first mass}
\label{arithdetail:first-mass}
The relative-class estimate of Appendix~\ref{app:history-details} is
\[
 \sum_h e^{-\delta C_y(h)}(q_h^{\rm rel})^\delta
 \leq C\Big(\prod_jD_j^{\rm own}\Big)^\delta.
\]
The exact outputs give
$R_y=\sum_jy_j+\epsilon/r_m-E_y$, with
$E_y=\sum_{j<m}d_j\eta_j^y\geq0$ and $-C\leq\epsilon\leq0$.
Since $\prod_jD_j^{\rm own}=e^{-\sum_jy_j}$,
$e^{R_y}\prod_jD_j^{\rm own}\leq1$.
Also, with the actual class $Q_h(Y)$,
\[
 \Phi_g(h)\sum_{Z\in Q_h(Y)}F(Z)=e^{I_h}q_h^{\rm rel},
 \qquad C_y(h)=R_y-I_h.
\]
Multiplication proves the unnormalized row bound
\[
 \sum_h\Phi_g(h)^\delta
       \Big(\sum_{Z\in Q_h(Y)}F(Z)\Big)^\delta\leq C
 \qquad(F(Y)>0).
\]
Every original pivot is retained. A bounded clock discrepancy changes
the constant once per original cut.

On harmonic tuple measure times Lebesgue measure let
$f(\mathbf a,t)=\sum_YF(Y)1_{B_Y(\mathbf a)}(t)$ and
$Z_F=\sum_YF(Y)$. The weight is bounded, invariant under permutations
inside each cell and independent of the actual prime values there.
By \eqref{arithdetail:uniform-supply},
$\int f=(\log2)^kS_gZ_F\geq cW_gZ_F$.
Pointwise, if $t\in B_Y$, every other box containing $t$ has
center difference less than $\log2$ in each coordinate. Consequently
\[
 f^{1+\delta}
 \leq\sum_YF(Y)1_{B_Y}
       \Big(\sum_{Z:Y\sim_{\mathbf p}Z}F(Z)\Big)^\delta.
\]
Integrate and apply Theorem~\ref{thm:correlated-collision-kernel} and
the row bound to obtain $\int f^{1+\delta}\leq CW_gZ_F$.
H\"older on its support gives
\[
 \operatorname{meas}\{f>0\}
 \geq\frac{(\int f)^{(1+\delta)/\delta}}
             {(\int f^{1+\delta})^{1/\delta}}
 \geq cW_gZ_F.
\]
Its support is contained in the original divisor union. Summing over
the disjoint occupancies retains exactly the first power of the
selected nominal mass.

\subsection{Bounded perturbations with the physical optimizer fixed}
\label{arithdetail:frozen-grid}
Fix the geometric, scalar, count, REG and occupancy constants, then
choose the roughness and cell grids. The initial roughness deletion
has bounded harmonic mass, and omitting the final incomplete cell
loses less than one fixed cell mass in each group. The coarse prime
sum comparison  therefore bounds the total resulting
length loss by a constant $R=O_k(1)$. Choose the lift above all source thresholds and above
$C_H(1+R)^2$. At the lifted physical vector $L$ optimize once and keep
its products, owner partition, fine probabilities, selected edges and
ordinary suffix fixed. The complete-cell lengths $L'$ satisfy
$\sum_i|L_i'-L_i|\leq R$. Every original group remains present.

Map each grid block affinely onto its physical block. This preserves
the order of activations, deadlines and eligible partitions of every
lawful history. All reference integrands are bounded on the compact
parameter ranges. Thus, uniformly over the full history space,
\begin{align}
 |\overline C_{L',s}(h')-\overline C_{L,s}(h)|&\leq C_HR,
 &\overline C_{L',s}(h')&\geq-C_HR,\notag\\
 |\partial_{s_i}J(L',s)-\partial_{s_i}J(L,s)|&\leq C_HR,
 &|J(L',s)-J(L,s)|&\leq C_HR.
 \label{arithdetail:bounded-defects}
\end{align}
The fixed-coefficient reference charges for wrong names, old weak
directions, mixed times, pure lifetimes and terminal rank curvature
acquire just this bounded additive defect. For example the same two
lawful competitors give
$\overline C_{\rm float}\geq c_HT_{\rm float}-C_HR$.
The old-mixture node equation has residual $O_H(R)$: divide by the
old length to obtain its entropy error and multiply back when making
the old-group charge. The final error is additive $O_H(R)$.
A nonempty old collection is large after the lift, so the hierarchy
$V_{\rm old}\leq C_HT_{\rm owner}+C_HR$ retains its fixed-factor
form. Compression and interpolation are performed on $L'$ itself;
only their KKT lower bounds have this intercept. The exact likelihood,
projection and carry identities have no stationarity error.

Two endpoint rows require exact treatment. For one owner let
$m_{\rm ref}=\sum_g(\mathbb EH_g-r_gL_g')=O_H(R)$ and
$\zeta_g=H_g-\mathbb EH_g$. Its accepted exact output is
$Y=a_0-b_0=m_{\rm ref}+\sum_g\zeta_g$, so
\begin{align}
 C_v&=a_0-\sum_{g\leq l}\zeta_g+\overline C_v'
                         +\sum_{\rm proper}\nu_{v,g},\notag\\
 &=b_0+m_{\rm ref}+\sum_{g>l}\zeta_g+\overline C_v'
                         +\sum_{\rm proper}\nu_{v,g}.
 \label{arithdetail:exact-grid-endpoints}
\end{align}
For the all-zero template $C_v=a_0$ exactly. For the complete-deadline
template $\nu=0$ and $\overline C_v'=-m_{\rm ref}$, so $C_v=b_0$
exactly and its intrinsic cost is zero. Neither small endpoint is
enlarged by $R$.
Every other endpoint has an own-scale charge. If its last positive
time is strictly internal, clipped charging gives at least
$c_H\min_gL_g$. If it reaches the end but is not everywhere full,
some proper group has either a linear bad-group charge or the displaced
good-group square-root charge. Hence
\begin{equation}
 \overline C_v+S_{\rm good}(v)
 \geq c_H\sqrt{\min_gL_g}-C_H.
 \label{arithdetail:nonexceptional-charge}
\end{equation}
The chosen lift absorbs all bounded defects within the already
reserved fractions of these charges. The displaced named box is
retained; central uncolored counts alone would not imply this step.

The finite context sets at $L,L'$ have identical indexing. Formal
perturbed scores may use $\max(0,\overline C_v')$, with the two
exceptional templates given their exact scores. Reference errors
are $O_H(R)$ and
$|\sqrt{L_g'}-\sqrt{L_g}|\leq\sqrt{|L_g'-L_g|}$.
Native exit scales are comparable. If an early/late flag changes,
$|t_v-T/2|\leq C_HR$. For large $T$ its reference charge is
$\geq c_HT-C_HR$, dominating both possible additions; for bounded
$T$, the fixed positive score floor suffices. Thus each paired score,
and then its separate minimum, changes by at most a fixed factor.
The actual grid pair still uses
$\Delta_i'=c_iL_i'-a_in_i+b_ik_i$ in the max construction.

At natural mean $\mu_i'=Z_iL_i'$, the centered parameter
$p_i^0=(a_i-c_i/Z_i)/b_i$ is unchanged; the gap-clock parameter
$L_i'/(\mu_i'+1)$ changes by $O_H(R)/(1+L_i)^2$.
The local Lipschitz estimate for $\chi$ and
$|\log T_i|\leq C_H\log(2+L_i)$ on the rare branch therefore
give a bounded exponent-times-log error. Score comparability gives
comparability of all binary factors, including threshold crossings.
Similarly $w_j'=R_j'/\sqrt{T_j'}\asymp w_j$.
The frozen ordinary suffix lengths change by $O_H(R)$; its numerical
factor is stable with $1+\max(D',0)$ in this comparison. If two
largest blocks are comparable, choosing either places the other on
a side and gives a fixed positive ballot factor. This numerical
observation does not itself admit a negative mean into a probability
theorem; the exact root argument below handles that possibility.

\subsection{Count bands and the exact root on the final grid}
\label{arithdetail:grid-counts}
For the same complete integer type, all grouped outputs and all
recursive exit coordinates change by deterministic $O_H(R)$ because
the count coefficients are fixed. Condition on the common pair-coarsened
field in the strict-exit construction. Each proper macro split remains
an independent $\operatorname{Bin}(m_j,p_j)$, with $m_j\asymp T_j$,
compact $p_j$, and the same coefficient in its exit; its effect on every
deeper exit cancels exactly. Band centers move by $O_H(R)$ without
reading any remaining split.

Translate each split by a deterministic integer of size $O_H(R)$,
starting in a fixed interior subband. The fixed width floor exceeds
the finitely many lattice steps, so rounding stays inside the original
band with a fixed fractional margin. The small change in its width
is absorbed by that margin. For a bounded translation the logarithm
of the binomial probability ratio on the central range is
$O_H(R/\sqrt{m_j}+R^2/m_j)$, directly by factorial ratios.
The central windows have spare width since $\sqrt{m_j}\gg R$.
The same factorial calculation applies to each finite multinomial
thick box: its conditional center and incoming count shift by
$O_H(R)$, while its fixed $\sqrt{L_g}$ width is unchanged.
This check is made stagewise before any path is sampled.

The natural uncolored means also change by $O_H(R)$.
For $|n-\mu|\leq C\sqrt\mu$, the logarithm of their Poisson ratio is
$n\log(\mu'/\mu)-(\mu'-\mu)=O_H(1)$.
Hence integration of the complete count measure retains its true
weights and mass $W'\asymp\prod_{j<m}w_j$.
Equations \eqref{arithdetail:exact-grid-endpoints} and
\eqref{arithdetail:nonexceptional-charge} supply the same strong-box
fractions and legal physical-score endpoint boxes. Terminal completion
is checked by the same thick-box calculation at every admitted root
count, retaining the original root clock.

The actual root count is the floor solving the exact grid equation.
At the same other counts it changes by at most $C_H(1+R)$ from the
physical-budget floor. Since $V_{\rm pre}\leq C_HU$ and
$U\gg(1+R)^2$, the Euler identity with its bounded residual gives
$|n_*'-\mu_*'|\leq C_H\sqrt U$. Its Poisson atom is comparable to
$U^{-1/2}$, uniformly over retained prefix environments and the
allowed logarithmic targets. All $n_*'+1$ gaps are then sampled at
this count.

Let $P$ be the entire original length before the ordinary weak block
and $B$ its later positive-drift length. If $P>0$, the lift makes it
large; the physical relation $D\asymp P$ gives
$D_{\rm grid}\geq c_HP-O_H(R)>0$. The actual ordinary pair remains
in the boxes $X\asymp1+P$, $Y\asymp1+B$, so the ordinary exact-type
theorem applies. If $P=0$, there are no earlier owners or terminal
binary head, all common targets vanish, and $x=0$. The root still
sets the total risk to its bounded floor error even if the natural
grid mean has either sign $O_H(R)$. If $B>0$, the later central
blocks have actual surplus at least
$c_HB-C_H\sqrt B-O_H(R)\geq c_H'B$. The weak first-block increment
is exactly determined by those counts and the root floor: one legal
endpoint remains fixed and the other is comparable to $1+B$.
If $B=0$, their difference is the bounded root floor and both
reserves remain fixed.

Apply the all-reserve ordinary exact-type strong-wall and REG theorem
directly to these legal pairs. Its hypotheses are the exact endpoint
relation, positive reserves and compact marks, speed and gap parameter;
they do not require a nonnegative natural mean. Other suffix blocks
retain their separated signed drifts at central counts. Conditional
on all counts, the original clocks and their constant-retention named
refinements give one joint kernel
\[
 \asymp U^{-1/2}
 \min\left\{1,\frac{(1+P)(1+B)}{1+M_{\rm weak}}\right\}
\]
when a weak block exists, and the saturated path factor otherwise.
The physical positive-tail identities identify this with the frozen
ordinary factor. Both sides use the same central root atom and the
same full-gap ordinary bridge; a later largest positive block makes
the path factor saturated.

All binary blocks now have supported central types, their actual grid
increments and positive common boxes. The curved-wall, REG and
all-transfer estimates apply at the fixed physical products. Integrate
the mixed moments at a fixed full type and then sum the complete type
measure. The exact owner powers and tilt identity give
\[
 K_k(\log3+L)\geq
 c_He^{-J(L';s)}\mathcal P_{\rm fixed}(L';s)
 \geq c_He^{-J(L;s)}\mathcal P_I(L).
\]
The first inequality is the first-mass transfer already proved; the
second uses \eqref{arithdetail:bounded-defects} and the factor
comparisons above. Products have remained fixed at the physical
optimizer throughout. All REG and occupancy constants remain those
used to choose the grid; only the final length threshold and additive
geometric constants depend on its loss bound. Finally
Theorem~\ref{thm:ell-stability} compares the original vector with its
one fixed lift, as used in Theorem~\ref{lower:lifted}.
